\documentclass{article}
\usepackage{graphicx} % Required for inserting images
\usepackage[utf8]{inputenc}
\usepackage[T1]{fontenc}
\usepackage{url}

\title{Magisterka}
\author{268784 }
\date{April 2026}

\begin{document}

%\maketitle

\section{Spotkanie 8.04}
Pocz\k{a}tkowa teoria: wg pracy (Starck, Jean-Luc, Emmanuel J. Cand\`es, David L. Donoho. "The curvelet transform for image denoising.")\\
Warto zrozumie\'c jak rozr\'o\.znia\'c FFT, $\hat{f}$, DFT, wavelet/curvelet/ridgelet transform oraz transformatę Radona.\\
Dobrym krokiem do dok\l{}adniejszego zrozumienia tematu b\k{e}dzie zaimplementowanie podstawowego modelu.\\
Seminarium 16.04: warto pokaza\'c og\'olniejszy zarys problemu. Jak si\k{e} da, to nie ogranicza\'c si\k{e} do jednego \'zr\'od\l{}a.
\section{Spotkanie 15.04}
Pocz\k{a}tek: 11:15\\
Curvelet transform, wavelet transform: Jean-Luc Starck, Fionn Murtagh, Albert Bijaoui. "Image processing and data analysis The multiscale approach" (książka z wprowadzeniem wielu pojęć , które wykorzystuje Starck)\\\\
Syllabus kursu Mathematical Image Processing (literatura na str. 3):\\ \url{https://wmat.pwr.edu.pl/d/FGBUKOTtQKxVvBEJ0ShNuClBQSXp5XUFABWdbExYbWGFBRAg9RgNtcXccVjU3ClVBDkNFeXVeW1cU/mathematical_image_processing.pdf}\\\\
Laura Freijeiro-Gonz\'alez, Manuel Febrero-Bande, Wenceslao
Gonz\'alez-Manteiga. "A critical review of LASSO and its derivatives for variable
selection under dependence among covariates" (przegląd metod regresji penalizowanej, w tym metod podobnych do LASSO):\\
\url{https://arxiv.org/pdf/2012.11470}\\\\
16.04 15-16, podpis ws. specjalności mgr
\section{Spotkanie 6.05}
Warto zacząć pisanie artykułu. W tym celu można na temat każdego znalezionego źródła napisać krótką notatkę dla siebie (w dowolnym miejscu), żeby móc później z jej skorzystać i wiedzieć, które źródło zacytować w danym momencie.\\\\
MR Lookup (dobra strona do wyszukiwania cytowań):\\
\url{https://mathscinet.ams.org/mrlookup}\\\\
\noindent
Przejście do metod regresji liniowej.\\
Warto sprawdzić, jak w praktyce redukuje się odczyt fMRI do macierzy (n obserwacji na p czynników).\\
Propozycja: każdy wycinek (slice) mózgu do odszumienia osobno [do sprawdzenia w źródłach dot. fMRI na ile to dobre i jeśli nie to czym to zastępują].\\\\

\noindent
Źródła do wykorzystania:
\begin{itemize}
    \item Alan Agresti. "Foundations 
of Linear and Generalized  Linear Models"
\item Trevor Hastie, Robert Tibshirani, Jerome Friedman. "The Elements of Statistical Learning" \url{https://web.stanford.edu/~hastie/ElemStatLearn/}
\item Źródła o odszumianiu obrazu w ogólnym przypadku, nie tylko fMRI
\end{itemize}
Wspomniane metody:
\begin{itemize}
    \item GLM: różne przypadki możliwe do wykorzystania i porównania
    \item ICA
    \item LASSO oraz Basis Pursuit
    \item filtr Kalmana (lepiej skupić się na poprzednich metodach)
\end{itemize}

\section{Spotkanie 11.06 (15.06)}
\noindent
Spotkanie w czwartek od 16:00 do 17:00, pokój A.4.6.\\
Przeniesiono na 15.06 (zdalnie)

\section{Spotkanie 17.06}
\noindent
Spotkanie w środę od 12:00 do 13:00, pokój A.4.6\\\\

\noindent
Następne spotkanie: piątek (10.07.2026), 11:00, pokój A.4.6


\end{document}
